% Pillar-7 conclusion
\section{Conclusion}
\label{sec:p7-conclusion}
The companion Pillar-2 paper ended by refusing to over-read a diagnostic;
this paper earns the reading it withheld. A one-line binomial model,
$S = p(1-p)(1 - h_G(p))$, predicts the ZVF phenomena Pillar~2 catalogued:
its exact expected value (T1, matched at the endpoints and interior of the
binned reward tensors), its monotone decline in group size (T2, confirmed
on two independent grids), and its coupling to gradient mass through
$p(1-p)$ (T3, directionally confirmed in an open backward pass at
$+0.71$). The same model exposes the diagnostic's structural limit---the
symmetry that aliases mastery with incapacity, visible as a U-shape across
368 runs---and supplies the repair: the pairwise-contrast density
$\mathrm{PCD} = \frac{G-1}{G}\E[p(1-p)]$, which survives the micro-jitter
that silently zeroes ZVF in any pipeline with a dense sub-reward.

Prediction enabled intervention. The \texttt{zvf-triage} callback and
adaptive-$G$ controller act on T2 live, escalating $G$ on starvation spikes;
in a four-arm open audit the controller matched the best fixed recipe's
held-out gain while dynamic sampling bought perfect telemetry at $+45\%$
rollouts with no outcome advantage. The resulting design rules are modest
and, we believe, durable: gate on magnitude (PCD), alarm on existence
(ZVF), escalate multiplicatively on spikes, price every intervention in
rollouts, and never let a recipe label substitute for a stack report. What
remains is the decisive tier---the cross-anchor PCD horse race, the
hard-population controller test, and production-scale gradient
attribution---which we have specified precisely so that this paper, like
its companion, can be falsified rather than merely believed.
